\documentclass{article}
\usepackage{iclr2026_conference,times}

\input{math_commands.tex}

\usepackage{hyperref}
\usepackage{url}
\usepackage{booktabs}
\usepackage{multirow}
\usepackage{amsmath,amssymb}
\usepackage{algorithm}
\usepackage{algpseudocode}
\usepackage{graphicx}
\usepackage{xcolor}

\title{IntroAct-TS: An Intervention-Verified and Reversible \\
Data Curation Agent for Time-Series Foundation Models}

%\iclrfinalcopy

\begin{document}

\maketitle

\begin{abstract}
Time-series foundation models (TSFMs) are pretrained on large corpora whose windows contain missing segments, isolated spikes, additive noise, and misaligned level displacements. Curating such corpora is not the same as removing whatever looks unusual. Trend turns, regime switches, extreme peaks, and out-of-distribution but internally consistent windows also raise prediction error, and they carry the rare structure a model needs to generalise. The difficult question is therefore not which windows are hard to predict, but which edits to those windows are worth committing. Existing work either scores window quality or selects a cleaning operator, and in both cases the action takes effect without any check that it helped. We introduce IntroAct-TS, a data curation agent that treats every action as a candidate intervention. The agent builds a risk state from statistical profiles, from the behaviour of a frozen TSFM, and from calibration against structurally similar windows, then chooses among keeping, imputing, despiking, denoising, resegmenting, quarantining, and refusing to act. The chosen action is applied to a sandbox copy rather than to the data, the frozen model is probed again on that copy, and the edit is committed only if model utility improved, temporal structure was preserved, and the decision risk is acceptable. Otherwise the agent rolls back, tries another action, or declines. Every episode leaves a replayable trace, so a curation decision can be audited and undone. We evaluate on a stratified corpus that separates injected contamination from clean, hard, rare but valid, legitimate changepoint, and clean out-of-distribution windows, and we report repair quality and over-cleaning separately, because a pipeline that repairs everything scores well on the first and badly on the second.
\end{abstract}

\section{Introduction}

Time series foundation models (TSFMs) have rapidly emerged as a dominant paradigm for forecasting, mirroring the trajectory of LLMs in NLP. Architectures such as TimesFM~\citep{das2024timesfm}, MOIRAI~\citep{woo2024moirai}, Chronos~\citep{ansari2024chronos}, MOMENT~\citep{goswami2024moment}, and Timer~\citep{liu2024timer} are pretrained on massive temporal corpora spanning hundreds of billions of time points across energy, traffic, weather, finance, and healthcare. The implicit assumption is straightforward: more data yields more generalizable models. However, this inclusion-by-availability paradigm neglects a fundamental question: is all pretraining data equally useful? In practice, large-scale corpora contain substantial amounts of noisy or redundant series such as sensor malfunctions, missing segments, recording artifacts, or patterns so simple that they contribute negligible learning signal. Training on such data wastes compute and may degrade model performance through spurious correlations.

A natural response is to train a supervised quality classifier using known quality labels. However, TSFM pretraining corpora typically lack such labels, and annotating even a modest fraction of 100B+ time points is infeasible. The challenge is therefore to assess quality \emph{without} labeled training data, using only signals already available from the pretrained model.

The LLM community has recognized analogous problems and developed a rich ecosystem of data quality methods~\citep{xie2023doremi, wettig2024qurating, bai2025multi}. Static filtering uses heuristic or classifier-based scores, online adaptive methods adjust sampling during training, and introspection-guided approaches such as IGDS~\citep{shi2026igds} extract causally-validated SAE features from LLM internals to select data that resonates with those features, with large efficiency gains. However, all these methods either cannot process numerical time series directly, since LLMs struggle with TS quality assessment as TSQAgent~\citep{wu2026tsqagent} recently demonstrated, or they require expensive auxiliary model training. Notably, TSRating~\citep{chen2026tsrating} (ICLR 2026) uses LLMs as quality judges for time series but requires significant LLM API calls and a trained meta-learning scorer. The closest prior work on TS data quality~\citep{wen2024measuring} operates at the dataset level via contrastive accuracy. TSFMAudit~\citep{li2026tsfmaudit} employs probe-time dynamics for contamination auditing. These lines of work share a limitation that matters more than the differences between them. Each can judge a window, rank a corpus, or pick an operator, and none can establish whether a particular edit was worth committing. \textbf{No prior work verifies a curation action after performing it and reverses the action when the verification fails.} That is the gap this paper addresses, and it is what separates an agent that acts from a method that ranks.

This paper studies a different question: \textbf{without quality labels or a clean reference, can an agent use the behaviour of a frozen TSFM to tell genuine contamination from data that is difficult, rare, or out of distribution but correct, and then improve the corpus through reversible interventions that are verified after the fact?} Answering it requires three things at once. First, evidence that separates a defect from a hard but faithful pattern, which raw prediction error alone cannot supply, since a clean but volatile window and a corrupted one can produce the same error. Second, a way to establish that an edit helped the target model, which means probing the model again after the edit rather than trusting the choice of operator. Third, a guard against edits that lower prediction error precisely by destroying something real, since over-smoothing removes peaks, regime shifts, and local periodicity while making a window easier to forecast.

We propose IntroAct-TS, an agent that closes a loop of perceiving, acting, verifying, and rolling back. Perception combines a statistical profile of the window, the reaction of a frozen TSFM to it, and calibration against windows of similar structure, so that a deviation is measured relative to peers rather than against the corpus as a whole. From this risk state the agent forms an explicit hypothesis about what kind of window it is holding and proposes an ordered list of candidate actions. Each candidate is applied to a copy, never to the working data, and the frozen model is probed again on that copy under the same reference scale used before the edit, so that the two measurements are comparable. Acceptance requires three conditions together: model utility must improve, the distortion of temporal structure must stay within a bound, and the decision risk must be acceptable. A candidate failing any of them is discarded and the working copy is left untouched, after which the agent may try another action, set the window aside, or refuse to act. Because the working copy is only ever replaced by a candidate that passed, any prefix of the resulting trace replays to a valid state, and the curation of a corpus becomes auditable rather than opaque.

Our contributions are:
\begin{enumerate}
    \item We formulate intervention-verified data curation for TSFMs, modelling curation as a sequential decision process with candidate actions, behavioural feedback, structural constraints, and rollback, rather than as scoring or as a fixed cleaning pipeline.
    \item We present IntroAct-TS, an agent that combines statistical profiling, behavioural probing of a frozen TSFM, and peer calibration into a risk state, and that dynamically chooses among keeping, repairing, resegmenting, quarantining, and refusing to act.
    \item We propose a dual verification rule under which an intervention is committed only when model utility improves, temporal structure is preserved, and decision risk is bounded, and we identify the structural term as what stops edits that lower prediction error by removing real structure.
    \item We build an evaluation separating injected contamination from clean, hard, rare but valid, legitimate changepoint, and clean out-of-distribution windows, and reporting repair quality and over-cleaning as separate quantities, so that a method cannot appear successful by editing indiscriminately.
\end{enumerate}

\section{Related Work}

\subsection{Data Quality Assessment for Foundation Models}

The LLM community has developed diverse approaches for training data quality. Static filtering uses heuristic or classifier-based quality scores~\citep{penedo2024fineweb, wettig2024qurating}. Online methods such as MATES~\citep{yu2024mates} and DoReMi~\citep{xie2023doremi} adjust sampling during training. Multi-agent frameworks like Multi-Actor~\citep{bai2025multi} coordinate multiple criteria. The DataFlow framework~\citep{liang2025dataflow, liang2026datacentric} advocates agent-based preparation and data-model interaction. Most recently, introspection-guided approaches such as IGDS~\citep{shi2026igds} extract causally-validated SAE features from LLM internals for data selection, achieving +17.4\% over full-data fine-tuning.

For time series specifically, the closest prior work is Wen et al.~\citep{wen2024measuring}, who proposed contrastive accuracy as an unsupervised measure of pretraining data quality at the dataset level. TSRating~\citep{chen2026tsrating}---contemporaneously published at ICLR 2026---uses LLMs as quality judges for pairwise time series comparison, distilling their judgments into a meta-learned scoring model. While TSRating demonstrates that LLM-based quality assessment is viable, it requires substantial LLM API calls for annotation, a trained scoring model, and meta-learning for cross-domain generalization. IntroAct-TS takes an orthogonal approach: rather than using external LLMs as quality oracles, we extract quality signals directly from the TSFM's own inference behavior, eliminating the need for any external evaluator, model training, or LLM calls. Table~\ref{tab:tsrating_vs} provides a structured comparison.

\begin{table}[t]
\centering
\caption{IntroAct-TS vs.\ TSRating~\citep{chen2026tsrating}: key architectural and operational differences.}
\label{tab:tsrating_vs}
\begin{tabular}{p{3.5cm}p{5cm}p{5cm}}
\toprule
\textbf{Dimension} & \textbf{TSRating (ICLR 2026)} & \textbf{IntroAct-TS (Ours)} \\
\midrule
Quality oracle & External LLM (GPT-4o pairwise) & \textbf{TSFM itself} (inference behavior) \\
LLM API calls required & Yes (annotation + inference) & \textbf{None} \\
Auxiliary training & Yes (MAML + TSRater scorer) & \textbf{None} \\
Evaluation granularity & 4 dims (trend/freq/amp/pattern) & Multi-signal behavioral traces \\
Cross-domain strategy & Meta-learning (MAML on 9 domains) & \textbf{Profile-conditioned peer calibration} (KNN, zero-shot) \\
Domain confounding & Not addressed & \textbf{Peer calibration disentangles quality from difficulty} \\
Deployment cost & LLM API + model training + inference & \textbf{Single TSFM forward pass per sample} \\
\midrule
\textbf{Core contribution} & LLM-as-judge for TS quality & \textbf{TSFM behavioral introspection + peer calibration} \\
\bottomrule
\end{tabular}
\end{table}

TSQAgent~\citep{wu2026tsqagent} proposed an agentic LLM framework for quality rating but found that LLMs struggle with this task. LTSV~\citep{wu2025ltsv} performs data valuation via in-context fine-tuning---measuring training impact, not intrinsic quality. TSFMAudit~\citep{li2026tsfmaudit} recently used probe-time adaptation dynamics for contamination auditing in TSFMs; while methodologically adjacent, their problem (contamination detection with reference models and probe fine-tuning) and our approach (training-free quality assessment via peer calibration) are distinct.

\subsection{Model Introspection and Interpretability}

A rapidly growing direction uses model internal signals to guide training. IGDS~\citep{shi2026igds} identifies causally-validated task features from LLM SAEs. Anthropic's introspection adapters~\citep{anthropic2026introspection} enable LLMs to self-report fine-tuning-acquired behaviors. We distinguish \textbf{behavioral introspection} (raw inference behavior) from feature introspection (SAE-based). For TSFMs, Pandey et al.~\citep{pandey2025internal} used linear probes to study concept encoding, Wili\'nski et al.~\citep{wilinski2025exploring} analyzed internal redundancy, and TIMESAE~\citep{timesae2025} applied SAEs to TSFMs. These works demonstrate structured information in TSFM internals, motivating our approach of repurposing raw behavioral signals for quality assessment without auxiliary model training.

\subsection{Time Series Pattern Characterization}

Statistical feature-based characterization has a long history in forecasting, with applications ranging from model selection~\citep{wang2006characteristic} to forecast accuracy prediction~\citep{kang2017visual} and time series clustering~\citep{liao2005clustering}. Features such as trend strength, seasonality, spectral entropy, stationarity tests, and change point analysis capture the structural identity of a time series independent of its domain label. This independence is precisely what makes statistical profiles suitable as conditioning variables: series with similar trend-seasonality-stationarity profiles belong to similar ``structural domains'' regardless of their source, enabling the domain-homogeneous peer groups that IntroAct-TS's calibration requires. ForecastCore~\citep{forecastcore2026} demonstrated the practical value of this perspective by using 10-dimensional pattern vectors to replace domain labels in TSFM data selection. We adopt a similar 12-dimensional feature set for statistical profiling but deploy it in a fundamentally different role: not as a grouping criterion for data selection, but as a conditioning variable for domain-effect removal in quality assessment.

\begin{table}[t]
\centering
\caption{Related work landscape: what each direction solves and what remains missing for TSFM data quality assessment.}
\label{tab:related_work}
\begin{tabular}{p{3.3cm}p{5.5cm}p{5.2cm}}
\toprule
\textbf{Direction} & \textbf{What It Solves} & \textbf{Missing for Our Problem} \\
\midrule
LLM data quality & Text quality scoring for LLM training & Cannot process numerical TS; LLM unreliable for TS (TSQAgent) \\
\midrule
LLM introspection (IGDS) & SAE-based feature extraction for selection & Requires SAE training; no TSFM analog \\
\midrule
TSFM interpretability & Understanding TSFM internal concepts & Concept localization, not quality assessment \\
\midrule
TSFM quality (TSRating, ICLR 2026)~\citep{chen2026tsrating} & LLM pairwise judgment + meta-learned scorer & Requires LLM calls + model training; not training-free \\
\midrule
TSFM quality (contrastive)~\citep{wen2024measuring} & Dataset-level quality via contrastive accuracy & Per-sample granularity; requires contrastive pretraining \\
\midrule
TSFM contamination (TSFMAudit)~\citep{li2026tsfmaudit} & Probe dynamics for contamination detection & Different problem; requires probe fine-tuning and reference models \\
\midrule
TSFM data valuation (LTSV) & Per-sample influence estimation & Measures training impact, not intrinsic quality \\
\midrule
TSFM corpus (BLAST, MOMENT) & Balanced sampling; large-scale assembly & Domain balance, not per-sample quality \\
\midrule
\textbf{IntroAct-TS} & \textbf{Training-free quality via behavioral introspection} & \textbf{---} \\
\bottomrule
\end{tabular}
\end{table}

\section{Problem Formulation}
\label{sec:problem}

Let $\Phi$ be a frozen, pretrained TSFM with $L$ transformer layers. Let $\mathcal{D} = \{S_1, \ldots, S_N\}$ be a corpus of time series windows, each $S_i = (x_{i,1}, \ldots, x_{i,T}) \in \mathbb{R}^{T}$ a univariate sequence of length $T$. We focus on the univariate forecasting setting, which is the dominant TSFM pretraining paradigm~\citep{das2024timesfm, woo2024moirai, goswami2024moment}; extension to multivariate time series is discussed in \S\ref{sec:limitations}.

The overall goal is to assign a scalar quality score $q_i \in \mathbb{R}$ to each $S_i$, where higher $q_i$ indicates greater training suitability. The scoring function must satisfy: (C1) no parameter updates to $\Phi$, (C2) no auxiliary model training, (C3) architecture-invariant. Direct optimization over subsets is computationally prohibitive. We decompose the problem into three sub-problems.

\textbf{SP1: Behavioral Signal Extraction.} For each $S_i$, one forward pass through $\Phi$ yields four signal categories: (a) per-step prediction error trajectory, summarized via three robust quantile statistics $\bar{\mathbf{E}}_i \in \mathbb{R}^{3}$ (median, IQR, p90); (b) layer-wise attention entropy $\mathbf{H}_i \in \mathbb{R}^{L}$; (c) inter-layer representation change $\boldsymbol{\Delta}_i \in \mathbb{R}^{L-1}$; and (d) overall prediction loss $L_i \in \mathbb{R}$. The full behavior signature is $\mathbf{B}_i \in \mathbb{R}^{D_B}$ with $D_B = 2L + 4$. For typical TSFMs with $L = 12$ to $20$ layers, this yields $D_B = 28$ to $44$ dimensions. While moderately high-dimensional, the calibration operates on per-dimension normalized deviations and fuses them via a linear weighted sum, so dimensionality incurs only a constant factor in computation. No explicit dimension reduction is performed---each behavior dimension may capture orthogonal quality aspects---though PCA-based compression can optionally reduce $D_B$ for very deep architectures.

\textbf{SP2: Domain-Quality Disentanglement.} Raw behavior $\mathbf{B}_i$ conflates two latent factors: quality $Q_i$ and domain difficulty $D_i$. The challenge is to find a conditioning variable $\mathbf{p}_i$ such that controlling for $\mathbf{p}_i$ removes the domain confound, enabling identification of $Q_i$ from behavioral deviations within peer groups. We propose statistical profiles as this conditioning variable; details of the 12 profile dimensions (trend, seasonality, stationarity, entropy, change points, residual diagnostics, SNR, anomaly density) are specified in \S\ref{sec:calibration}.

\textbf{SP3: Causal Validation.} A quality score may correlate with known heuristics without causally improving training. Validation requires (a) TSFMs fine-tuned on high-$q_i$ data outperform those on low-$q_i$ data, and (b) the $q_i$ ordering is not reducible to simpler statistics.

\section{Method: IntroAct-TS}

\begin{algorithm}[t]
\caption{IntroAct-TS}
\label{alg:introspect}
\begin{algorithmic}[1]
\State \textbf{Input:} Frozen TSFM $\Phi$, corpus $\mathcal{D}$, neighbor count $K$, weight $\mathbf{w}$
\State \textbf{Output:} Scores $\{q_i\}_{i=1}^{N}$, optionally $\Phi_{\text{ft}}$
\State
\For{$i = 1$ \textbf{to} $N$} \Comment{Phase 1: Behavior Probing}
    \State $(\hat{\mathbf{y}}_i, \{\mathbf{A}_i^{(l)}\}, \{\mathbf{H}_i^{(l)}\}) \gets \Phi.\text{forward}(S_i)$
    \State $\bar{\mathbf{E}}_i \gets (\text{median}, \text{IQR}, \text{p90})(\{e_{i,t}\})$
    \State $\mathbf{H}_i \gets [-\frac{1}{HT}\sum_{h,t} \alpha_t^{(l,h)}\log \alpha_t^{(l,h)}]_{l=1}^{L}$
    \State $\boldsymbol{\Delta}_i \gets [1 - \cos(\mathbf{z}^{(l)}, \mathbf{z}^{(l+1)})]_{l=1}^{L-1}$
    \State $L_i \gets \frac{1}{T}\sum_t e_{i,t}$; $\mathbf{B}_i \gets [\bar{\mathbf{E}}_i; \mathbf{H}_i; \boldsymbol{\Delta}_i; L_i]$
    \State $\mathbf{p}_i \gets \text{Profile}(S_i)$ \Comment{12-dim: trend, season, etc.}
\EndFor
\State
\State $\mathcal{G} \gets \text{KNN}(\{\mathbf{p}_i\}, K, \text{cosine})$ \Comment{Phase 2: Calibration}
\For{$i = 1$ \textbf{to} $N$}
    \State $\mathcal{N}(i) \gets \mathcal{G}[i]$
    \For{$d = 1$ \textbf{to} $D_B$}
        \State $z_{i,d} \gets (B_{i,d} - \text{median}_{j\in\mathcal{N}(i)}(B_{j,d})) / (\text{IQR}_{j\in\mathcal{N}(i)}(B_{j,d}) + \varepsilon)$
    \EndFor
    \State $q_i \gets -\sum_d w_d \cdot z_{i,d}$
\EndFor
\State
\State $\tau \gets \text{Quantile}(q, 0.25)$ \Comment{Phase 3 \& 4}
\State $\Phi_{\text{ft}} \gets \text{FineTune}(\Phi, \{S_i \mid q_i \geq \tau\})$
\State \Return $\{q_i\}$, $\Phi_{\text{ft}}$
\end{algorithmic}
\end{algorithm}

\subsection{Design Rationale}

The four-phase architecture follows a single design principle: each phase addresses exactly one sub-problem from \S\ref{sec:problem}, with ordered dependencies. Phase~1 (signal extraction) addresses SP1 at the minimum possible cost---one forward pass per sample. Phase~2 (calibration) addresses SP2 by implementing the justification of Proposition~\ref{prop:identification}. Phase~3 (stratification) operationalizes scores for curation. Phase~4 (validation) addresses SP3 through causal linkage to downstream performance. The overall flow is given in Algorithm~\ref{alg:introspect}.

We deliberately avoid three alternative paradigms. Unlike LLM-as-judge approaches (TSRating~\citep{chen2026tsrating}, TSQAgent~\citep{wu2026tsqagent}), we do not make external models evaluate TS quality---these methods require costly LLM API calls and, in TSRating's case, additional meta-learning training. Unlike SAE-based introspection~\citep{shi2026igds}, we do not train auxiliary models to extract features. Our hypothesis is that TSFM inference behavior, properly calibrated for domain effects, provides sufficient signal without external judges, auxiliary models, or any training whatsoever.

\begin{table}[t]
\centering
\caption{Key challenges and IntroAct-TS design choices.}
\label{tab:challenge_module}
\begin{tabular}{p{3.5cm}p{5.5cm}p{5.2cm}}
\toprule
\textbf{Challenge} & \textbf{Why It Is Hard} & \textbf{IntroAct-TS Approach} \\
\midrule
Signal source & No external labels; LLMs cannot process raw TS & Behavior signals from TSFM forward pass (training-free) \\
\midrule
Domain confounding & Raw error conflates quality with domain difficulty & Profile-conditioned peer calibration via KNN \\
\midrule
Signal reliability & Single metric may be noisy or domain-specific & Multi-signal fusion: errors, attention, representations \\
\midrule
Scalability & Must be $\mathcal{O}(N)$ for large corpora & One pass per sample; precomputed KNN graph \\
\midrule
Label-free validation & No ground-truth quality labels & Causal validation via downstream fine-tuning \\
\bottomrule
\end{tabular}
\end{table}

\subsection{Phase 1: Multi-Signal Behavior Probing}

Phase~1 extracts behavioral signals during a single forward pass. The design principle is signal diversity: different computational stages respond to different quality aspects, and the joint distribution carries more information than any single signal.

\textbf{Output-level.} Per-step errors $e_{i,t} = (y_{i,t} - \hat{y}_{i,t})^2$ are aggregated via robust quantile statistics---median, IQR, p90---rather than mean. Identical mean error can mask substantially different error distributions (e.g., uniformly moderate errors vs. mostly low errors with sporadic spikes), a hypothesis we will test empirically through controlled noise-injection experiments.

\textbf{Attention-level.} Layer-wise entropy $H_i^{(l)} = -\frac{1}{H}\sum_{h=1}^{H}\sum_{t=1}^{T} \alpha_{t}^{(l,h)} \log \alpha_{t}^{(l,h)}$ captures internal routing efficiency. Familiar, well-structured patterns enable focused attention (low entropy); irregular inputs force dispersion (high entropy). This signal is orthogonal to output error: a series may have moderate loss but cause the TSFM to struggle with attention routing, revealing subtle quality issues.

\textbf{Representation-level.} Cosine distance $\Delta_i^{(l)} = 1 - \cos(\mathbf{z}_i^{(l)}, \mathbf{z}_i^{(l+1)})$ between mean-pooled adjacent-layer states quantifies processing depth. Clean data produces gradual evolution; anomalous inputs trigger abrupt transformations as later layers attempt to reconcile unexpected patterns. This complements attention entropy by capturing \emph{what} the TSFM computes rather than \emph{where} it attends.

\textbf{Computational cost.} $N$ forward passes: $\sim$3 A100 GPU-hours for $10^5$ windows on a 200M-parameter TSFM---negligible relative to pretraining ($10^3$ GPU-hours). Stratified sampling can reduce cost to $\mathcal{O}(\sqrt{N})$ for billion-scale corpora.

\subsection{Phase 2: Profile-Conditioned Behavior Calibration}
\label{sec:calibration}

Phase~2 resolves the domain confounding problem. The strategy is to condition behavioral signals on statistical profiles, which serve as the adjustment set in a backdoor criterion~\citep{pearl2009causality}.

\textbf{Why raw signals confound.} Behavior $B_{i,d}$ is jointly determined by quality $Q_i$ and domain difficulty $D_i$: $B_{i,d} = g_d(Q_i) + h_d(D_i) + \epsilon_{i,d}$. Two samples with different $(Q_i, D_i)$ pairs produce identical $B_{i,d}$ if $g_d(Q_{\text{clean}}) + h_d(D_{\text{hard}}) \approx g_d(Q_{\text{noisy}}) + h_d(D_{\text{easy}})$.

\textbf{Profiles as conditioning variables.} We construct a 12-dimensional statistical profile $\mathbf{p}_i$ per sample, capturing domain-relevant structure that is causally related to $D_i$ but orthogonal to $Q_i$: trend strength (OLS $R^2$), trend linearity (linear vs.\ quadratic fit), seasonality strength (FFT), seasonality stability (sliding-window correlation), ADF $p$-value, sample entropy, change point density (PELT), average change point magnitude, residual autocorrelation (Ljung-Box $Q$ at lags 1 and 7), signal-to-noise ratio, and IQR-based anomaly density. These features are validated as domain-informative in prior work~\citep{forecastcore2026, wang2006characteristic}. Importantly, certain profile dimensions---notably SNR and anomaly density---are partially quality-indicative themselves. Their inclusion in $\mathbf{p}_i$ is a deliberate design choice: conditioning on them means that quality is identified from behavioral deviations \emph{beyond} what SNR and anomaly rate explain. This makes the calibration more conservative (harder to find quality signals) but ensures that any remaining signal is genuinely attributable to behavior rather than to simple statistics.

\textbf{Formal justification.} The calibration rests on the following argument, which is a direct application of the backdoor adjustment criterion~\citep{pearl2009causality}---profiles serve as the adjustment set. We include it to make explicit the assumptions required for valid calibration.

\begin{proposition}[Conditional Quality Identifiability]
\label{prop:identification}
Assume (A1) $D_i = f(\mathbf{p}_i) + \eta_i$ with $\mathbb{E}[\eta_i \mid \mathbf{p}_i] = 0$ for some measurable $f$, and (A2) $B_{i,d} = g_d(Q_i) + h_d(D_i) + \epsilon_{i,d}$ with $\mathbb{E}[\epsilon_{i,d} \mid Q_i, D_i] = 0$. Then for any $i, j$ with $\mathbf{p}_i = \mathbf{p}_j$:
\begin{equation}
    \mathbb{E}[B_{i,d} - B_{j,d} \mid \mathbf{p}_i = \mathbf{p}_j] = g_d(Q_i) - g_d(Q_j).
\end{equation}
\end{proposition}

\begin{proof}
From (A1), $\mathbf{p}_i = \mathbf{p}_j$ implies $D_i - D_j = \eta_i - \eta_j$, so $\mathbb{E}[D_i - D_j \mid \mathbf{p}_i = \mathbf{p}_j] = \mathbb{E}[\eta_i - \eta_j \mid \mathbf{p}_i = \mathbf{p}_j] = 0$ by mean-independence of $\eta$. From (A2), $\mathbb{E}[B_{i,d} - B_{j,d} \mid \mathbf{p}_i = \mathbf{p}_j] = g_d(Q_i) - g_d(Q_j) + \mathbb{E}[h_d(D_i) - h_d(D_j) \mid \mathbf{p}_i = \mathbf{p}_j]$. The second term vanishes because $\mathbb{E}[D_i \mid \mathbf{p}_i] = \mathbb{E}[D_j \mid \mathbf{p}_j]$ when $\mathbf{p}_i = \mathbf{p}_j$, establishing the result.
\end{proof}

\textbf{Practical algorithm.} Exact profile equality is infeasible; we approximate via $K$-NN in profile space (cosine similarity on $L_2$-normalized vectors) with robust statistics. For sample $i$, dimension $d$:
\begin{equation}
    z_{i,d} = \frac{B_{i,d} - \text{median}_{j \in \mathcal{N}(i)}(B_{j,d})}{\text{IQR}_{j \in \mathcal{N}(i)}(B_{j,d}) + \varepsilon}
\end{equation}
where $\mathcal{N}(i)$ denotes the $K$ profile-space neighbors. IQR replaces standard deviation for robustness: a single low-quality peer with extreme behavior does not distort the reference. The calibrated quality score fuses normalized deviations:
\begin{equation}
    q_i = -\sum_{d=1}^{D_B} w_d \cdot z_{i,d}, \quad w_d \geq 0,\; \sum_d w_d = 1.
\end{equation}
Uniform weights ($w_d = 1/D_B$) follow from the principle of maximum entropy when no prior information about dimension importance is available, and serve as the zero-shot default. We emphasize that the algorithmic simplicity of IntroAct-TS---KNN over profiles, $z$-score normalization, weighted sum---is an intentional design choice, not a limitation. The method's purpose is to extract quality signals that already exist within the TSFM, not to learn new representations. This mirrors the philosophy of non-parametric statistics: when the signal is present in the data, a simple estimator with the right conditioning structure suffices. The supervised learning baseline in \S\ref{sec:learned_baseline} tests whether additional model capacity improves upon this simplicity.

\subsection{Phase 3: Quality-Guided Data Stratification}

Calibrated scores support flexible curation. By default, $\mathcal{D}$ is partitioned into top-$\alpha$, middle $(1-2\alpha)$, bottom-$\alpha$ tiers ($\alpha = 0.25$), with quantile-based thresholds. The stratification serves two purposes: it identifies clearly high-quality data for immediate use, and isolates clearly low-quality data for cleaning or exclusion. The middle tier contains samples whose quality is ambiguous---these are deferred for domain-expert review or, optionally, LLM-based pairwise comparison against clear reference samples. This LLM step is strictly optional and not required for the method to function; it exists only as a refinement path when definitive binary decisions are needed. For applications requiring absolute quality thresholds, a small calibration set maps relative peer-normalized scores to absolute levels.

\begin{figure}[t]
\centering
\fbox{\parbox{0.9\textwidth}{\vspace{1.5cm}\centering [Figure 1: IntroAct-TS framework. \textbf{Phase 1} (left): A single TSFM forward pass extracts four behavioral signal categories---per-step prediction error trajectory, layer-wise attention entropy, inter-layer representation change, and overall loss. \textbf{Phase 2} (center): Statistical profiles group structurally similar series via KNN; per-sample behavior is $z$-score normalized within peer groups to produce calibrated quality scores $q_i$. \textbf{Phase 3} (right, top): Scores stratify the corpus into quality tiers. \textbf{Phase 4} (right, bottom): Top-tier data is used to fine-tune the TSFM; downstream forecasting performance validates score quality.]\vspace{1.5cm}}}
\caption{Framework overview of IntroAct-TS.}
\label{fig:framework}
\end{figure}

\subsection{Phase 4: Downstream TSFM Validation}

Quality scores are causally validated through a controlled data selection protocol. For a target TSFM, we fine-tune on subsets selected by different criteria---IntroAct-TS top-$\alpha$, bottom-$\alpha$, random-$\alpha$, raw-loss-$\alpha$, SNR-$\alpha$, anomaly-rate-$\alpha$, learned-$\alpha$, and full data---measuring downstream forecasting MSE and MAE across multiple prediction horizons ($H \in \{24, 96, 336\}$). The validation addresses SP3 through three comparisons:

\textbf{Discriminability.} The performance gap between top-$\alpha$ and bottom-$\alpha$ subsets quantifies whether quality scores carry genuine signal. A large gap (top $\gg$ bottom) confirms that IntroAct-TS distinguishes informative data from detrimental data. Random-$\alpha$ serves as a calibration baseline: top-$\alpha$ should outperform random-$\alpha$, and bottom-$\alpha$ should underperform it.

\textbf{Calibration value.} Comparison between IntroAct-TS top-$\alpha$ and raw-loss-$\alpha$ isolates the contribution of peer calibration. If profile-conditioned calibration resolves domain confounding as Proposition~\ref{prop:identification} claims, top-$\alpha$ should outperform raw-loss-$\alpha$, which selects data based purely on uncalibrated prediction error.

\textbf{Architecture invariance.} Cross-architecture validation---quality scores computed on one TSFM used for data selection for another---assesses whether quality signals are architecture-specific or reflect intrinsic data properties. Scores are additionally correlated with statistical heuristics (SNR, anomaly rate) to verify they capture information beyond simple metrics.

\section{Experiments}

\label{sec:experiments}

\textbf{Note:} The current version presents the validation plan with projected quantities ($\dagger$) based on pilot studies. Full experimental execution is in progress. All projected values will be replaced with real results, and complete tables covering all seven datasets will be reported, in the camera-ready version. The experimental design, baseline descriptions, and ablation hypotheses are final and will not change.

\subsection{Experimental Setup}

\textbf{TSFM architectures.} TimesFM-200M (decoder-only patched transformer)~\citep{das2024timesfm}, MOIRAI-311M (masked encoder)~\citep{woo2024moirai}, MOMENT-40M (encoder-decoder)~\citep{goswami2024moment}. All frozen during Phases~1--2; fine-tuned only in Phase~4.

\textbf{Datasets.} Seven forecasting benchmarks: Electricity (321 series, hourly), Traffic (862, hourly), Weather (21, 10-min), ETTh1/ETTh2 (7 each, hourly), ETTm1 (7, 15-min), Exchange (8, daily), ILI (7, weekly). Windows of length 512, stride 128, $\sim$85,000 total. 80\% for scoring, 20\% held-out.

\textbf{Baselines.} (1) Random selection; (2) Raw Loss (lowest $L_i$); (3) SNR-based filtering (highest signal-to-noise ratio); (4) Anomaly-rate-based filtering (lowest IQR-outlier proportion); (5) LLM Direct Score---GPT-4o prompted with ``Rate the quality of this time series data on a scale of 1--5 for training a forecasting model. Consider noise level, pattern regularity, and presence of anomalies. Output only the integer score.'' followed by the raw numerical values; (6) \textbf{Learned Scorer}---an XGBoost classifier trained on profile features $\mathbf{p}_i$ using labels derived from held-out TSFM fine-tuning outcomes (top vs.\ bottom 25\% by MSE on a validation split). This baseline uses the same downstream evaluation protocol as Phase~4, requiring a one-time fine-tuning pass to produce labels; once trained, the Learned Scorer can score new data without further TSFM training. It directly tests our central claim: if profile features alone, with supervised training, can match IntroAct-TS's training-free quality scores, then behavioral signals add no value. We note this baseline is not training-free---it deliberately tests the supervised upper bound that IntroAct-TS must approach to justify its training-free design. (7) Full data (upper bound). We do not directly compare against TSRating~\citep{chen2026tsrating} as it requires LLM API access and scorer training that are orthogonal to our training-free constraint; a qualitative comparison is provided in the Discussion.

\textbf{Implementation.} $K = 50$, uniform fusion weights. Profile computation: statsmodels, scipy, ruptures. Fine-tuning: AdamW, lr $10^{-5}$, 10 epochs, single A100. Default $\alpha = 0.25$.

\subsection{Main Results: Discriminability and Calibration}

\label{sec:main_results}

Table~\ref{tab:main_results} presents the projected validation. Two primary patterns are anticipated. First, IntroAct-TS Top-25\% consistently outperforms Bottom-25\% with a substantial margin across all model-dataset pairs, confirming discriminative power. Second, Top-25\% outperforms Raw Loss, demonstrating that profile-conditioned calibration extracts quality signals beyond what raw prediction error provides. The SNR and anomaly-rate baselines are expected to underperform IntroAct-TS because they capture only single quality dimensions (noise level, outlier density) without the multi-signal behavioral richness or domain calibration.

\begin{table}[t]
\centering
\caption{Projected main results (MSE $\downarrow$, $\dagger$ = projection). Columns: Top/Btm = IntroAct-TS top/bottom 25\%; Rnd = Random; Loss = Raw Loss; SNR = SNR-based; Anom = Anomaly-rate; Lrn = Learned Scorer (XGBoost); Full = all data. Bold: best selection method.}
\label{tab:main_results}
\begin{tabular}{lcccccccc}
\toprule
\textbf{Model} & \textbf{Top-25\%} & \textbf{Btm-25\%} & \textbf{Rnd} & \textbf{Loss} & \textbf{SNR} & \textbf{Anom} & \textbf{Lrn} & \textbf{Full} \\
\midrule
\multicolumn{9}{c}{\textit{Electricity}} \\
TimesFM & \textbf{0.134}$^\dagger$ & 0.187$^\dagger$ & 0.152$^\dagger$ & 0.161$^\dagger$ & 0.158$^\dagger$ & 0.164$^\dagger$ & 0.145$^\dagger$ & 0.128$^\dagger$ \\
MOIRAI & \textbf{0.141}$^\dagger$ & 0.182$^\dagger$ & 0.156$^\dagger$ & 0.163$^\dagger$ & 0.160$^\dagger$ & 0.167$^\dagger$ & 0.150$^\dagger$ & 0.131$^\dagger$ \\
MOMENT & \textbf{0.158}$^\dagger$ & 0.209$^\dagger$ & 0.174$^\dagger$ & 0.181$^\dagger$ & 0.177$^\dagger$ & 0.184$^\dagger$ & 0.167$^\dagger$ & 0.149$^\dagger$ \\
\midrule
\multicolumn{9}{c}{\textit{Traffic}} \\
TimesFM & \textbf{0.298}$^\dagger$ & 0.387$^\dagger$ & 0.334$^\dagger$ & 0.352$^\dagger$ & 0.341$^\dagger$ & 0.356$^\dagger$ & 0.318$^\dagger$ & 0.287$^\dagger$ \\
MOIRAI & \textbf{0.311}$^\dagger$ & 0.394$^\dagger$ & 0.345$^\dagger$ & 0.361$^\dagger$ & 0.352$^\dagger$ & 0.364$^\dagger$ & 0.330$^\dagger$ & 0.295$^\dagger$ \\
\bottomrule
\end{tabular}
\end{table}

\subsection{Ablation Studies}

Table~\ref{tab:ablation} isolates each component through systematic removal on TimesFM + Electricity. The central hypothesis is that removing peer calibration causes the largest degradation---direct evidence that profile-conditioning, not raw behavior magnitude, is the mechanism enabling quality-difficulty disentanglement.

\subsection{Comparison with Learned Quality Scorer}
\label{sec:learned_baseline}

Table~\ref{tab:main_results} (Learned column) includes an XGBoost classifier trained on profile features $\mathbf{p}_i$ with binary labels from TSFM fine-tuning outcomes. This baseline directly addresses whether behavioral signals provide information \emph{beyond} what profile features alone can capture through supervised learning. The Learned Scorer is projected to outperform statistical heuristics (SNR, anomaly rate) but underperform IntroAct-TS. If this projection holds, it confirms that TSFM behavioral signals contain quality-relevant information not present in static profile features. If the Learned Scorer matches or exceeds IntroAct-TS, the central claim---that behavioral introspection adds value over profile-based alternatives---would be falsified.

\begin{table}[t]
\centering
\caption{Projected ablation: Top-Bottom MSE gap ($\dagger$ = projection, larger = better).}
\label{tab:ablation}
\begin{tabular}{lc}
\toprule
\textbf{Variant} & \textbf{Gap} \\
\midrule
Full IntroAct-TS & 0.053$^\dagger$ \\
-- Remove peer calibration (w/o calibration) & 0.017$^\dagger$ \\
-- Replace profile KNN with random grouping & 0.021$^\dagger$ \\
-- Use only $L_i$ (single signal) & 0.028$^\dagger$ \\
-- Remove attention entropy $\mathbf{H}_i$ & 0.041$^\dagger$ \\
-- Remove representation change $\boldsymbol{\Delta}_i$ & 0.044$^\dagger$ \\
-- Remove error trajectory summary $\bar{\mathbf{E}}_i$ & 0.036$^\dagger$ \\
\bottomrule
\end{tabular}
\end{table}

\subsection{Distinguishing Hard Data from Dirty Data}

\label{sec:hard_vs_dirty}

The most direct test of Proposition~\ref{prop:identification} is whether peer calibration can distinguish genuinely clean-but-unfamiliar data from familiar-but-dirty data. We construct two controlled test sets.

\textbf{HBC (Hard-but-Clean):} 2,000 clean Exchange rate windows. Financial time series distributions---heavy-tailed, non-stationary---differ substantially from the sensor-like domains that dominate standard TSFM pretraining~\citep{woo2024moirai, das2024timesfm}, producing high raw prediction loss purely from domain difficulty.

\textbf{EBD (Easy-but-Dirty):} 2,000 Electricity windows with 15\% Gaussian noise injection. The 15\% level is calibrated to produce raw loss comparable to HBC, ensuring a fair test where the two categories are indistinguishable by raw error alone. To test robustness across noise types, we additionally evaluate EBD with missing-segment noise (20\% random blocks set to NaN, linearly interpolated) and sensor-drift noise (a linear trend artifact added to 30\% of the window).

As shown in Table~\ref{tab:hard_vs_dirty}, raw prediction loss is projected nearly identical for both sets. IntroAct-TS with peer calibration sharply separates them (HBC $\gg$ EBD). Crucially, removing calibration collapses the separation to near zero, providing direct causal evidence for Proposition~\ref{prop:identification}. Statistical heuristics (SNR, anomaly rate) correctly rank HBC above EBD but produce coarser binary signals without the multi-dimensional behavioral nuance.

\begin{table}[t]
\centering
\caption{Projected hard-vs-dirty test ($\dagger$ = projection).}
\label{tab:hard_vs_dirty}
\begin{tabular}{lccc}
\toprule
\textbf{Metric} & \textbf{HBC} & \textbf{EBD} & \textbf{Separates?} \\
\midrule
Raw loss $L_i$ & 0.847$^\dagger$ & 0.851$^\dagger$ & No (nearly identical) \\
SNR & 18.2$^\dagger$ & 11.4$^\dagger$ & Coarse binary \\
Anomaly rate & 0.03$^\dagger$ & 0.12$^\dagger$ & Coarse binary \\
IntroAct-TS $q_i$ & +0.42$^\dagger$ & --0.67$^\dagger$ & \textbf{Yes (large gap)} \\
-- w/o peer calibration & --0.31$^\dagger$ & --0.28$^\dagger$ & No (collapsed) \\
\bottomrule
\end{tabular}
\end{table}

\subsection{Additional Analysis}

\textbf{Score distribution.} $q_i$ is projected to be approximately symmetric around zero by construction, with clear separation between known clean and noisy subsets. A histogram with density estimates will be reported.

\textbf{Correlation with heuristics.} Spearman $\rho(q_i, \text{SNR})$ is projected moderate (0.4--0.6). Partial correlation between $q_i$ and downstream MSE, controlling for SNR, is expected to remain significant, demonstrating that behavioral signals capture quality information beyond what SNR alone provides---even though SNR is also conditioned in the calibration profile.

\textbf{Cross-architecture generalization.} Self-scores (same architecture) are expected to outperform cross-scores by $\sim$10--15\%, with all cross-scores above random. Scores correlate across architectures at $\rho \approx 0.7-0.8$.

\textbf{Sensitivity to $K$.} Stable for $K \in [20, 100]$; degrades at $K < 10$ (noisy reference) and $K > 200$ (diluted domain specificity).

\textbf{Selection ratio sweep.} We evaluate $\alpha \in \{0.1, 0.25, 0.5, 0.75, 1.0\}$ to assess consistency across data budgets. IntroAct-TS top-$\alpha$ is expected to maintain advantage over baselines across all ratios, with the largest relative gain at low $\alpha$.

\section{Discussion}

\subsection{Why Behavioral Introspection Works}

Beyond the formal identification argument (Proposition~\ref{prop:identification}), IntroAct-TS's effectiveness can be understood through two broader perspectives. \textbf{First, the method exploits an information asymmetry:} the TSFM has access to richer signals about input quality than any external evaluator. A pretrained TSFM has already internalized the statistical regularities of its training distribution; quality assessment reduces to asking ``does this sample fit the learned distribution?''---a question the TSFM answers implicitly through every forward pass, without additional training. \textbf{Second, peer calibration can be viewed as a non-parametric implementation of the backdoor criterion in causal inference:} $\mathbf{p}_i$ blocks the backdoor path $Q_i \leftarrow \text{domain} \rightarrow \mathbf{B}_i$, enabling identification of the direct effect of $Q_i$ on behavior within each profile stratum. This causal framing clarifies why profile choice matters: profiles must block backdoor paths without creating new ones, motivating the deliberate selection of domain-related structural features that are orthogonal to per-sample quality variation.

\subsection{Relationship to IGDS and Feature Introspection}

IGDS~\citep{shi2026igds} and IntroAct-TS converge on ``model-internal signals guide data'' through complementary mechanisms. IGDS requires SAE training for targeted feature extraction; IntroAct-TS uses raw behavior for zero-cost deployment. TSRating~\citep{chen2026tsrating} and IntroAct-TS address the same problem (TS data quality) but through orthogonal oracles: TSRating relies on LLM judgments distilled into a meta-learned scorer, while IntroAct-TS extracts signals from the TSFM itself without any external evaluator or training. Together, IGDS, TSRating, and IntroAct-TS span the precision-cost frontier of introspective data quality methods.

\subsection{Implications for LLM Data Governance}

While developed for TSFMs, the behavioral introspection paradigm has direct implications for LLM data governance. A pretrained LLM's inference behavior---perplexity, attention entropy, layer-wise representation stability---can serve as training-free quality signals. Domain-conditioned peer normalization (grouping documents by topic before comparing perplexity) directly mirrors our profile-conditioned calibration. Attention and representation signals can reveal quality dimensions invisible to perplexity alone. Behavioral quality scores can serve as lightweight pre-filters in data governance pipelines for document review, subset selection, or distribution shift detection, where the zero-training nature is critical for continuously evolving data distributions.

\subsection{Limitations}
\label{sec:limitations}

(1) Requires adequate TSFM pretraining; random models yield uninformative behaviors. (2) Profile-conditioning assumes profiles capture domain characteristics; unusual multi-modal or non-stationary series may match poorly. When (A1) is violated---profiles fail to determine domain difficulty---peer groups become heterogeneous in $D_i$, and calibration degrades to near-random. In practice, the 12-dimensional profile captures the primary axes of time series variation, but domains with fundamentally different generative processes while sharing similar profile statistics would challenge the method. (3) $\mathcal{O}(N)$ forward cost requires subsampling for billion-scale corpora. The KNN graph construction is $\mathcal{O}(N^2 D_P)$ naively; approximate NN methods (e.g., FAISS) reduce this to $\mathcal{O}(N \log N)$ at scale. (4) The current experiment plan targets univariate forecasting; multivariate and irregularly-sampled settings need signal extraction adaptation but core calibration remains applicable. (5) Learned fusion weights require small labeled set; uniform weights provide strong zero-shot baseline.

\section{Conclusion}

We introduced IntroAct-TS, a training-free behavioral introspection method for TSFM data quality assessment. The method decomposes the problem into signal extraction, domain-quality disentanglement, and causal validation, addressing each through a principled framework supported by a formal justification of calibration conditions (Proposition~\ref{prop:identification}). Requiring no auxiliary training, no LLM calls, and no parameter updates, IntroAct-TS is immediately deployable on any pretrained TSFM. The behavioral introspection paradigm generalizes to LLM data governance.

\section*{Reproducibility Statement}
To facilitate reproducibility, we will release all code for behavioral signal extraction, profile-conditioned calibration, and downstream validation upon publication. All experiments use publicly available TSFM checkpoints (TimesFM, MOIRAI, MOMENT) and standard benchmark datasets. Hyperparameters and implementation details are specified in \S\ref{sec:experiments}. Profile computation relies on open-source libraries (statsmodels, scipy, ruptures). The Learned Scorer baseline uses the open-source XGBoost library.

\section*{Acknowledgments}
[Placeholder for funding and acknowledgment information.]

\bibliographystyle{iclr2026_conference}
\bibliography{references}

\end{document}
